\scanpage{599}

\mspage{8}

a le diagramme commutatif

\diagram{\begin{tikzcd}[ampersand replacement=\&,column sep=small]
\text{\unsure{(a)|(u)}}\quad\&k_E\arrow[r,dash,"\unsure{\chi|X}_\Phi","{\unsure{-|\to}}"{pos=1,anchor=center}]\arrow[rr,dash,bend right=24,"\chi_\Psi"',"{{\unsure{-|\to}}}"{pos=1,anchor=center}]\&k_\Phi\arrow[r]\&k_\Psi\ ,
\end{tikzcd}}

ce qui montre que $\Psi$ est dans l'image de \unsure{(a)|(u)}. Inversement, supposons que $\Psi\in\unsure{\mathscr F|\mathcal F}_0$ provienne de $\mathrm{Max}(k_\Phi)$, i.e. qu'\deleted{il existe $k_\Phi\to k_\Psi$}\inserted{l'on ait \deleted{\unsure{un}}} \deleted{\unsure{$f$}} $\unsure{\mathfrak m|m}_\Phi\subset\mathfrak m_\Psi$ \deleted{\unsure{(}} (\unsure{où|i.e.} $\mathfrak m_\Phi$ est le noyau de $\chi_\Phi$), ou encore qu'on ait une factorisation (a) de $\chi_\Psi$.

Prouvons que $\Psi$ est plus fin que $\Phi$, i.e. que $\Psi\supset\Phi$. Soit donc $A\in\Phi$, prouvons que l'on a $A\in\Psi$. \deleted{Cela va résulter du} \inserted{Mais $A\in\Phi$ \unsure{implique|signifie}} \deleted{Si on avait $A\notin\Psi$, on} fait que \deleted{\unsure{$f$}} $\chi_\Phi(\unsure{f_A|\varphi_A})=1$ \deleted{est \uline{inversible} dans $k_\Phi$,} \deleted{en conclurait que $f_A\in\mathfrak m_\Psi$} donc son image dans $k_\Psi$ est \deleted{\unsure{une unité|non nulle},} égale à 1, donc on a $A\in\Psi$ par le cor. 1. \deleted{\unsure{précédent}, Je \unsure{ferai}} \deleted{\unsure{donc} j'\unsure{utiliserai} ledit corollaire \unsure{ainsi} :}

\deleted{

% slanted
\margin{\unsure{remonté} à la place du cor. 1}

\uline{Corollaire 3} (de prop. 2) Soit $\Phi$ un filtre sur $\unsure{k|E}$. \deleted{Alors \unsure{$\Phi$} $A\in\Phi$} et $A\in\unsure{\mathcal P|\mathfrak P}(E)$. Alors

\[
A\in\Phi\Longleftrightarrow\chi_\Phi(f_A)\ \text{est inversible dans }k_\Phi
\]

(\unsure{donc|avec} $\Phi$ est \unsure{connu}, quand on \unsure{connait|connaît} l'anneau quotient $\chi_\Phi$ de $k_E$).
}

% archive sheet number: 598
